\begin{textsample}{NEG Dim 4 – 2019\_c4 – Score 2.00 – 2019\_c4/enade\_2019\_tec\_gest\_hosp\_50af38c3bf.txt}  \label{ex:f4_neg_246}
CAPÍTULO 4 PERCEPÇÃO DA PROVA Nas análises feitas neste capítulo, constam as percepções dos concluintes da Área de Tecnologia em Gestão Hospitalar sobre a prova aplicada no Enade/2019.
Essas percepções foram mensuradas por meio de nove questões que avaliaram desde o grau de dificuldade da prova até o tempo gasto para concluí -la.
As percepções sobre a prova foram relacionadas ao desempenho dos estudantes, à Grande Região de funcionamento do curso, à Categoria Administrativa e ao tipo de Organização Acadêmica da IES.
O questionário de percepção da prova encontra -se ao final do Anexo VII, que traz a reprodução do exame.
O desempenho dos estudantes foi classificado em quatro quartos.
Para tanto, esse desempenho foi ordenado de forma ascendente.
O percentil 25, P25, também conhecido como primeiro quartil, é a nota de desempenho que deixa um quarto ( 25 \% ) dos valores observados abaixo e três quartos acima.
A Figura 1 apresenta uma ilustração desse conceito.
O quarto inferior de desempenho é composto pelas notas abaixo do primeiro quartil.
Já o percentil 75, P75, também conhecido como terceiro quartil, é o valor para o qual há três quartos ( 75 \% ) dos dados abaixo e um quarto acima dele.
O quarto superior de desempenho é composto pelas notas iguais ou acima do terceiro quartil.
O percentil 50, P50, também conhecido como mediana, é o valor que divide as notas em dois conjuntos de igual tamanho.
O segundo quarto inclui valores entre o primeiro quartil ( P25 ) e a mediana.
O terceiro quarto contém os valores entre a mediana ( P50 ) e o terceiro quartil ( P75 ).
Vale ressaltar que percentis, quartis e medianas são pontos que, não obrigatoriamente, pertencem ao conjunto original de dados, ao passo que os quartos são subconjuntos dos dados originais.
A seguir, serão apresentados gráficos com resultados selecionados, relativos às nove questões avaliadas por grupos de estudantes.
Nas barras dos gráficos, apresenta -se o percentual de alunos que assinalaram uma das opções ou a soma das porcentagens daqueles que assinalaram duas ( ou três ) delas.
Por exemplo, para as questões 1 e 2, nos gráficos, é apresentada a porcentagem total de participantes que assinalaram as opções ( D ) Difícil e ( E ) Muito difícil.
Em cada barra, foram assinalados, também, os extremos do intervalo de confiança de 95 \% como linhas verticais unidas por uma linha horizontal na forma da letra H maiúscula.
O estimador de um parâmetro com um certo nível de confiança ( e.g. 95,0 \% ) deve conter o parâmetro no intervalo de confiança em 95 \% das vezes.
Na comparação entre os estimadores dos parâmetros de duas classes de uma dada categoria ( e.g.
Norte e Nordeste nas Grandes Regiões, de primeiro e último quarto dentro de desempenho, de IES Públicas e Privadas, ou de Universidades e Faculdades ) associados aos seus respectivos intervalos de confiança, diz -se que não há diferença estatisticamente significativa entre os parâmetros das duas categorias quando há interseção entre os intervalos de confiança e que há diferença, se os intervalos de confiança são disjuntos ( para maiores informações vide Glossário ).
Nas Tabelas no Anexo II, são apresentados os valores absolutos e a distribuição percentual18 das alternativas válidas das nove questões, segundo o mesmo recorte de desempenho dos alunos e Grande Região de funcionamento do curso, de Categoria Administrativa e do tipo de Organização Acadêmica da IES. 18 Cumpre lembrar uma das convenções para tabelas numéricas ( pág. iii ) sobre a possibilidade de a soma das partes não resultar em 100 \% por questões de arredondamento. 4.1 GRAU DE DIFICULDADE DA PROVA 4.1.1 Componente de Formação Geral Ao avaliarem Qual o grau de dificuldade desta prova na parte de Formação Geral? ( Questão 1 ), 39,3 \% do grupo de inscritos e presentes optaram pelas alternativas Difícil ou Muito difícil.
Entretanto, para mais da metade dos estudantes ( 55,9 \% ), o Componente de Formação Geral da prova foi considerado com grau de dificuldade Médio ( Gráfico 4.1, Gráfico 4.2 e, no Anexo II, a Tabela II.1 ).
O percentual de estudantes que consideraram a prova como Difícil ou Muito difícil foi maior na região Sudeste, onde a proporção foi de 41,9 \%, enquanto a de menor incidência foi a Nordeste, com 29,2 \%.
No Gráfico 4.1, é possível observar que a diferença entre a região Nordeste e a região Sudeste é estatisticamente significativa.
Nas Grandes Regiões, a proporção de presentes à prova que consideraram o Componente de Formação Geral como sendo de grau de dificuldade Médio esteve entre 53,3 \%, na região Sudeste e 66,1 \%, na região Nordeste.
O percentual de alunos que consideraram o componente de Formação Geral da prova Difícil ou Muito difícil foi decrescente em função dos quartos de desempenho, com exceção do primeiro para o segundo quarto : 40,8 \%, no primeiro quarto e 35,2 \%, no quarto, grupo de melhor desempenho na prova.
Nos quartos de desempenho intermediários, a proporção de alunos que consideraram a prova Difícil ou Muito difícil foi de 42,3 \%, no segundo quarto e de 39,0 \%, no terceiro quarto.
As diferenças entre todos os quartos de desempenho não são estatisticamente significativas.
Para todos os quartos de desempenho, a alternativa modal para essa pergunta foi a Médio, com 53,1 \% e 58,3 \% dos respondentes nos quartos extremos, primeiro e quarto, respectivamente.
O percentual de estudantes que consideraram o componente de Formação Geral da prova como Difícil ou Muito difícil foi maior para os estudantes de CEFET/IF, para os quais a proporção foi de 41,9 \%, enquanto as Faculdades foram a Organização Acadêmica com a menor incidência, com 34,5 \%, ainda que sem diferença estatisticamente significativa de Universidades, com 38,7 \%.
No Gráfico 4.3, é possível observar que as diferenças entre CEFET/IF, Faculdades, Centros Universitários ( 41,8 \% ) e Universidades não são estatisticamente significativas.
Nas Organizações Acadêmicas, a proporção de presentes à prova que consideraram o Componente de Formação Geral como sendo de grau de dificuldade Médio esteve entre 53,6 \%, nos Centros Universitários e 61,5 \%, nas Faculdades ( ver também a Tabela II.2, no Anexo II ).
O percentual de estudantes que consideraram o componente de Formação Geral da prova como Difícil ou Muito difícil foi maior para os estudantes de cursos de IES Privadas ( 39,7 \% ), sem uma diferença estatisticamente significativa da proporção em IES Públicas ( 36,0 \% ).
Já a proporção de presentes à prova que consideraram o Componente de Formação Geral como sendo de grau de dificuldade Médio foi 60,0 \% para os de IES Públicas e 55,4 \% para os de IES Privadas ( ver Gráfico 4.4 ). 4.1.2 Componente de Conhecimento Específico Ao responderem à Questão 2 – Qual o grau de dificuldade desta prova na parte de Componente Específico? – 34,7 \% do grupo de estudantes classificaram -na como Difícil ou Muito difícil.
Além disso, o Componente modal de Conhecimento Específico da prova foi considerado com grau de dificuldade Médio por 60,1 \% dos alunos ( Gráfico 4.5, Gráfico 4.6, e, no Anexo II, a Tabela II.3 ).
A análise das respostas dos estudantes quanto ao grau de dificuldade do Componente de Conhecimento Específico da prova, agregado por Grande Região, mostra que a diferença entre a menor e as duas maiores proporções de alunos que a avaliaram como Difícil ou Muito difícil é estatisticamente significativa, sendo a menor, na região Centro-Oeste ( 23,4 \% ), e as duas maiores, nas regiões Sudeste ( 37,0 \% ) e Sul ( 37,6 \% ).
A diferença entre a região Nordeste ( 25,6 \% ), a segunda menor proporção, e as regiões Sul e Sudeste, também é estatisticamente significativa.
O percentual de alunos que classificaram o grau de dificuldade como Médio, no Componente de Conhecimento Específico, variou de 58,2 \% a 70,2 \%, para as regiões Sul e Centro-Oeste, respectivamente.
Considerando -se a avaliação da dificuldade das questões do Componente de Conhecimento Específico da prova, de acordo com o desempenho dos estudantes, observa -se que há diferença estatisticamente significativa dos resultados entre o primeiro quarto e o quarto superior de desempenho.
A proporção dos que classificaram a parte específica como Difícil ou Muito difícil variou de 28,6 \% ( último quarto ) a 40,9 \% ( primeiro quarto ).
As demais proporções dessa reposta foram de 36,6 \% e 33,1 \%, nos quartos terceiro e segundo, respectivamente.
Já a proporção dos que responderam que o grau de dificuldade das questões do Componente de Conhecimento Específico da prova foi Médio variou de 52,9 \%, no primeiro quarto a 65,2 \%, no último.
Considerando -se a avaliação da dificuldade das questões do Componente de Conhecimento Específico da prova, o percentual de estudantes que as consideraram como Difícil ou Muito difícil foi maior para os estudantes de Universidades, para os quais a proporção foi de 37,4 \%, enquanto os CEFET/IF foram a Organização Acadêmica com a menor incidência, com 29,0 \%, ainda que sem diferença estatisticamente significativa de Faculdades, com 29,5 \%.
No Gráfico 4.7, é possível observar que as diferenças entre CEFET/IF, Faculdades, Centros Universitários ( 34,5 \% ) e Universidades não são estatisticamente significativas.
Nas Organizações Acadêmicas, a proporção de presentes à prova que consideraram o Componente de Formação Geral como sendo de grau de dificuldade Médio esteve entre 57,5 \%, nas Universidades e 64,5 \%, nos CEFET/IF ( ver também Tabela II.4, no Anexo II ).
O percentual de estudantes que consideraram o Componente Específico da prova como Difícil ou Muito difícil foi maior para os estudantes de cursos de IES Privadas ( 35,9 \% ), uma diferença estatisticamente significativa da proporção em IES Públicas ( 25,6 \% ).
Já a proporção de presentes à prova que consideraram este componente como sendo de grau de dificuldade Médio foi 65,6 \% para os de IES Públicas e 59,4 \% para os de IES Privadas ( ver Gráfico 4.8 e Tabela II.4, no Anexo II ). 4.2 EXTENSÃO DA PROVA EM RELAÇÃO AO TEMPO TOTAL Indagados quanto à extensão da prova, em relação ao tempo total oferecido para a sua resolução ( Questão 3 ), os estudantes apontaram, com maior incidência, a alternativa que considerava a extensão Adequada, para todas as agregações consideradas ( Gráfico 4.9, Gráfico 4.10, e, no Anexo II, a Tabela II.5 ).
O percentual de alunos que responderam ser a extensão da prova Adequada foi de 60,4 \%.
Já 33,7 \% dos inscritos presentes consideraram que a prova foi Longa ou Muito longa, e 5,9 \% a avaliaram como Curta ou Muito curta.
Dentre as Grandes Regiões, a proporção daqueles que avaliaram a prova como Longa ou Muito longa, em relação ao tempo total destinado à sua resolução, variou de 32,1 \%, na região Nordeste até 45,8 \%, na região Centro-Oeste.
Não há diferença estatisticamente significativa entre as regiões.
Considerando -se o desempenho dos alunos, nota -se ainda que 57,0 \% consideraram a extensão da prova Adequada, no quarto de desempenho inferior e 64,5 \%, no de melhor desempenho ( quarto superior ).
Nos quartos intermediários, essa proporção foi de 60,1 \%, no segundo e 59,6 \%, no terceiro quarto, proporções crescentes com o desempenho, exceto entre os quartos intermediários.
No Gráfico 4.6, pode -se constatar que há diferença estatisticamente significativa entre as proporções de estudantes que consideraram a prova Longa ou Muito longa, no primeiro quarto e no quarto de melhor desempenho, sendo a maior proporção dessa resposta, no primeiro quarto de desempenho ( 37,2 \% ), e a menor, no quarto superior ( 30,7 \% ), mas não caracterizando uma tendência.
Considerando -se a extensão da prova, o percentual de estudantes que avaliaram a prova como Longa ou Muito longa, em relação ao tempo total destinado à sua resolução, foi maior para os estudantes de Centros Universitários, para os quais a proporção foi de 35,7 \%, enquanto os CEFET/IF foram a Organização Acadêmica com a menor incidência, com 22,6 \%, com diferenças estatisticamente significativas de Universidades, com 34,5 \% e de Faculdades, com 31,0 \%.
No Gráfico 4.11, é possível observar que as diferenças de CEFET/IF em relação a Faculdades, Centros Universitários e Universidades são estatisticamente significativas.
Nas Organizações Acadêmicas, a proporção de presentes à prova que consideraram a sua extensão como adequada, em relação ao tempo total destinado à sua resolução, esteve entre 59,4 \%, nos Centros Universitários e 66,1 \%, nos CEFET/IF ( ver também Tabela II.6, no Anexo II ).
O percentual de estudantes que consideraram a prova como Longa ou Muito longa, em relação ao tempo total destinado à sua resolução foi maior para os estudantes de cursos de IES Privadas ( 34,6 \% ), sem uma diferença estatisticamente significativa da proporção em IES Públicas ( 26,4 \% ).
Já a proporção de presentes à prova que consideraram a extensão como Adequada foi 65,6 \% para os de IES Públicas e 59,7 \% para os de IES Privadas ( ver Gráfico 4.12 e Tabela II.6, no Anexo II ). 4.3 COMPREENSÃO DOS ENUNCIADOS DAS QUESTÕES 4.3.1 Componente de Formação Geral Em relação aos enunciados das questões do Componente de Formação Geral ( Questão 4 ), as opiniões foram positivas, já que 69,2 \% dos alunos avaliados consideraram todos ou a maioria dos enunciados das questões claros e objetivos ( Gráfico 4.13, Gráfico 4.14, e, no Anexo II, a Tabela II.7 ).
Na análise regional, a percentagem de estudantes que avaliaram que todos ou a maioria dos enunciados das questões do Componente de Formação Geral estavam claros e objetivos variou de 66,0 \%, na região Centro-Oeste a 75,4 \%, na região Norte.
A análise das percepções dos estudantes sobre a clareza e a objetividade dos enunciados permite afirmar que todos, ou a maioria dos enunciados de questões relativas ao Componente de Formação Geral, foram considerados claros e objetivos para a maior parte dos respondentes ( maior ou igual a 66,0 \%, em todas as regiões e maior ou igual a 61,1 \%, para todos os quartos de desempenho ).
Segundo o desempenho, observa -se que a proporção dos que emitiram essa opinião cresce conforme o desempenho aumenta, exceto do segundo para o terceiro quartos, com diferenças estatisticamente significativas entre o primeiro e os dois quartos de melhor desempenho.
No quarto superior, a clareza e a objetividade de todos ou da maioria dos enunciados das questões foram percebidas por 73,4 \%.
A proporção dos presentes que consideraram que todos ou a maioria dos enunciados das questões do Componente de Formação Geral estavam claros e objetivos foi maior para os estudantes de CEFET/IF, para os quais a proporção foi de 72,6 \%, enquanto as Faculdades foram a Organização Acadêmica com a menor incidência, com 68,0 \%, ainda que sem diferença estatisticamente significativa de Universidades, com 68,9 \% e de Centros Universitários, com 69,5 \%.
No Gráfico 4.15, é possível observar que as diferenças entre todas as organizações acadêmicas não são estatisticamente significativas ( ver também Tabela II.8, no Anexo II ).
O percentual de estudantes que consideraram que todos ou a maioria dos enunciados das questões do Componente de Formação Geral estavam claros e objetivos foi maior para os estudantes de cursos de IES Privadas ( 68,9 \% ), sem uma diferença estatisticamente significativa da proporção em IES Públicas ( 71,2 \% ).
Um maior detalhamento pode ser encontrado no Gráfico 4.16 e na Tabela II.8, no Anexo II. 4.3.2 Componente de Conhecimento Específico Em relação aos enunciados das questões do Componente de Conhecimento Específico da prova, para 77,7 \% dos estudantes avaliados da Área de Tecnologia em Gestão Hospitalar, a clareza e a objetividade ( Questão 5 ) estavam presentes em todas ou na maioria das questões ( Gráfico 4.17 a Gráfico 4.20, e no Anexo II, as Tabelas II.9 e II.10 ).
A maioria dos estudantes de todas as Grandes Regiões brasileiras considerou claros e objetivos todos ou a maioria dos enunciados das questões do Componente de Conhecimento Específico da prova, percentual sempre maior ou igual a 75,0 \%.
A diferença entre as regiões não é estatisticamente significativa.
A proporção de estudantes que consideraram os enunciados das questões claros e objetivos apresenta uma tendência crescente em relação ao aumento de desempenho, exceto entre o terceiro e o último quartos : mais elevada no terceiro quarto ( 85,0 \% ), se comparada ao quarto inferior de desempenho ( 67,9 \% ).
As diferenças do primeiro quarto de desempenho, em relação ao terceiro quarto e ao último, assim como a diferença entre o segundo quarto e o terceiro, são estatisticamente significativas.
A proporção dos presentes que consideraram que todos ou a maioria dos enunciados das questões do Componente de Conhecimento Específico estavam claros e objetivos foi maior para os estudantes de CEFET/IF, para os quais a proporção foi de 85,5 \%, enquanto os Centros Universitários foram a Organização Acadêmica com a menor incidência, com 76,4 \%, ainda que sem diferença estatisticamente significativa de Universidades, com 77,3 \% e de Faculdades, com 79,2 \%.
No Gráfico 4.19, é possível observar que as diferenças entre as organizações acadêmicas não são estatisticamente significativas ( ver também Tabela II.10, no Anexo II ).
O percentual de estudantes que consideraram que todos ou a maioria dos enunciados das questões do Componente de Conhecimento Específico estavam claros e objetivos foi maior para os estudantes de cursos de IES Públicas ( 85,6 \% ), sem uma diferença estatisticamente significativa da proporção em IES Privadas ( 76,7 \% ).
Ver Gráfico 4.20 e Tabela II.10, no Anexo II, para um maior detalhamento. 4.4 SUFICIÊNCIA DAS INFORMAÇÕES/INSTRUÇÕES FORNECIDAS Ao avaliarem as informações/instruções fornecidas para a resolução das questões ( Questão 6 ), 80,8 \% dos respondentes da Área de Tecnologia em Gestão Hospitalar de todo o Brasil afirmaram que estas eram até excessivas ou suficientes em todas ou na maioria das questões ( Gráfico 4.21, Gráfico 4.22, e, no Anexo II, a Tabela II.11 ).
Quanto à distribuição de respondentes pelas Grandes Regiões, observa -se que a proporção de estudantes que consideraram as informações/instruções fornecidas até excessivas ou suficientes em todas ou na maioria das questões foi sempre superior ou igual a 76,8 \%, chegando a 84,9 \%, na região Sudeste.
Não existem diferenças estatisticamente significativas entre as Grandes Regiões.
Levando -se em conta o desempenho dos participantes, observa -se uma tendência crescente das proporções de participantes que avaliaram as informações/instruções como até excessivas ou suficientes em todas ou na maioria das questões, havendo diferenças estatisticamente significativas dos dois quartos mais baixos, em relação aos dois quartos mais elevados de desempenho, como mostrado no Gráfico 4.12.
Nos dois quartos mais elevados, terceiro ( 85,7 \% ) e quarto superior ( 86,9 \% ), os percentuais foram acima do percentual da média nacional ( 80,8 \% ).
Já no quarto inferior, a suficiência das informações/instruções declarada como até excessiva, em todas ou na maioria das questões foi percebida por 74,4 \% dos respondentes e, no segundo quarto, por 76,2 \%.
Considerando -se as informações/instruções fornecidas para a resolução das questões ( Questão 6 ), a proporção que considerou que estas eram até excessivas ou suficientes em todas ou na maioria das questões foi maior para os estudantes de CEFET/IF, para os quais a proporção foi de 83,9 \%, enquanto as Faculdades foram a Organização Acadêmica com a menor incidência, com 76,4 \%.
No Gráfico 4.23, é possível observar que as diferenças das Faculdades em relação a CEFET/IF, Centros Universitários ( 81,7 \% ) e Universidades ( 81,4 \% ) não são estatisticamente significativas ( ver também Tabela II.12, no Anexo II ).
O percentual de estudantes que consideraram que as informações/instruções fornecidas para a resolução das questões eram até excessivas ou suficientes em todas ou na maioria das questões foi maior para os estudantes de cursos de IES Públicas ( 82,4 \% ), sem uma diferença estatisticamente significativa da proporção em IES Privadas ( 80,7 \% ) ( ver Gráfico 4.24 e Tabela II.12, no Anexo II ). 4.5 DIFICULDADE ENCONTRADA AO RESPONDER À PROVA Indagados sobre as dificuldades com as quais se depararam ao responder à prova ( Questão 7 ), 12,5 \% dos estudantes apontaram o Desconhecimento do conteúdo.
Para 51,1 \%, a Forma diferente de abordagem do conteúdo foi indicada como dificuldade.
Já a Falta de motivação para fazer a prova foi a dificuldade apontada por 9,9 \% dos respondentes.
Considerando -se todo o Brasil, 22,7 \% dos respondentes afirmaram que não tiveram qualquer tipo de dificuldade para responder à prova ( Tabelas II.13 e II.14, no Anexo II ).
Nos Gráficos 4.25 a 4.28, são apresentados os percentuais de estudantes que apontaram o Desconhecimento do conteúdo como dificuldade ao responder à prova.
Na análise por Grandes Regiões, o percentual de inscritos e presentes que apontaram o Desconhecimento do conteúdo como dificuldade ao responder à prova, não superou 16,7 \%.
Os percentuais variaram de 3,4 \%, na região Norte a 16,7 \%, na região Centro-Oeste.
A Forma diferente de abordagem do conteúdo foi a escolha modal dos estudantes, com percentuais que variaram de 47,6 \% ( região Nordeste ) a 56,9 \% ( região Norte ).
O percentual de alunos que citaram a Falta de motivação como dificuldade variou de 8,6 \% ( região Norte ) a 10,7 \% ( região Nordeste ).
Os que declararam não ter qualquer tipo de dificuldade para responder à prova variaram de 14,6 \%, na região Centro-Oeste a 29,3 \%, na região Norte.
Em relação aos quartos de desempenho, o Desconhecimento do conteúdo foi a opção escolhida por 10,7 \% dos estudantes do segundo quarto e por 13,8 \% do primeiro quarto, não caracterizando uma tendência.
Não são observadas diferenças estatisticamente significativas entre os quartos de desempenho.
A alternativa modal para os alunos, quando agregados pelos quartos de desempenho, para a dificuldade encontrada, foi causada pela Forma diferente de abordagem do conteúdo : 47,1 \%, no quarto superior e 55,1 \%, do primeiro quarto, escolheram essa alternativa.
Na análise por Organização Acadêmica, o percentual de inscritos e presentes que apontaram o Desconhecimento do conteúdo como dificuldade ao responder à prova não superou 16,1 \%.
Os percentuais variaram de 12,2 \%, para Universidades a 16,1 \%, nos CEFET/IF.
A Forma diferente de abordagem do conteúdo foi a escolha modal dos estudantes, com percentuais que variaram de 46,8 \% ( CEFET/IF ) a 53,5 \% ( Centros Universitários ).
O percentual de alunos que citaram a Falta de motivação como dificuldade variou de 7,5 \% ( Centros Universitários ) a 19,4 \% ( CEFET/IF ).
Os que declararam não ter qualquer tipo de dificuldade para responder à prova variaram de 14,5 \%, nos CEFET/IF a 27,6 \%, nas Faculdades.
No Gráfico 4.27, é possível observar que as diferenças entre CEFET/IF ( 16,1 \% ), Faculdades ( 12,6 \% ), Centros Universitários ( 12,4 \% ) e Universidades ( 12,2 \% ) não são estatisticamente significativas ( ver também Tabela II.14, no Anexo II ).
Em relação às Categorias Administrativas, o Desconhecimento do conteúdo foi a opção escolhida por 10,4 \% dos estudantes nas IES Públicas e por 12,8 \%, nas IES Privadas.
Não são observadas diferenças estatisticamente significativas entre estas duas categorias.
A alternativa modal para os alunos, quando agregados pelas Categorias Administrativas, para a dificuldade encontrada, foi causada pela Forma diferente de abordagem do conteúdo : 48,0 \%, para as Públicas e 51,4 \%, para as Privadas ( ver Gráfico 4.28 e Tabela II.14, no Anexo II ). 4.6 CONTEÚDOS DAS QUESTÕES OBJETIVAS DA PROVA Ao analisar os conteúdos das questões objetivas da prova ( Questão 8 ), um percentual muito pequeno dos estudantes avaliados, apenas 4,0 \%, afirmaram que não estudou ainda a maioria desses conteúdos ( Gráficos 4.29 a Gráfico 4.32, e as Tabelas II.15 e II.16, no Anexo II ).
A maioria ( 73,8 \% ) afirmou ter estudado e aprendido muitos ou todos os conteúdos avaliados.
Na análise por Grande Região, a proporção de respondentes que escolheram a opção não estudou ainda a maioria desses conteúdos foi pequena.
Observa -se que, nas regiões Norte ( 0,0 \% ), Nordeste ( 3,6 \% ) e Sudeste ( 3,5 \% ), as proporções foram menores do que a média nacional ( 4,0 \% ).
Nas regiões Sul ( 4,7 \% ) e Centro-Oeste ( 8,5 \% ), a situação foi inversa : uma proporção mais elevada.
Não são observadas diferenças estatisticamente significativas entre as regiões.
Em todas as regiões, a maioria dos presentes afirmou ter estudado e aprendido muitos ou todos os conteúdos, com proporções variando entre 68,1 \%, na região Centro-Oeste a 86,2 \%, na região Norte.
Considerando -se, separadamente, as opiniões de estudantes dos quatro quartos de desempenho, observa -se que, no quarto inferior, 8,4 \% ofereceram como resposta que não estudou ainda a maioria desses conteúdos, sendo 1,4 \% os do terceiro quarto com a mesma resposta.
A diferença entre os alunos que optaram por esse motivo de dificuldade no primeiro quarto e os dos dois quartos superiores de desempenho é estatisticamente significativa.
Tendo -se em conta o quarto superior, 83,7 \% dos alunos afirmaram ter estudado e aprendido muitos ou todos os conteúdos.
No outro extremo, no primeiro quarto, 56,8 \% optaram pelas mesmas categorias.
As proporções são crescentes com o desempenho.
Na análise por Organização Acadêmica, a proporção de respondentes que escolheram a opção não estudou ainda a maioria desses conteúdos foi pequena.
Observa -se que, nos Centros Universitários ( 3,6 \% ), Faculdades ( 3,4 \% ) e CEFET/IF ( 1,6 \% ), as proporções foram menores do que a média nacional ( 4,0 \% ).
Nas Universidades, a situação foi inversa : uma proporção mais elevada ( 4,7 \% ).
No Gráfico 4.31, é possível observar que as diferenças entre as organizações acadêmicas não são estatisticamente significativas ( ver também Tabela II.16, no Anexo II ).
Em todos os tipos de Organização Acadêmica, a maioria dos presentes afirmou ter estudado e aprendido muitos ou todos os conteúdos, com proporções variando entre 73,6 \%, nas Faculdades a 75,8 \%, nos CEFET/IF.
Em relação às Categorias Administrativas, observa -se que, nas IES Públicas, 0,8 \% ofereceu como resposta que não estudou ainda a maioria desses conteúdos, sendo 4,4 \% nas IES Privadas, com a mesma resposta.
A diferença entre os alunos que optaram por esse motivo de dificuldade nas duas categorias é estatisticamente significativa.
Considerando -se as IES Públicas, 84,8 \% dos alunos afirmaram ter estudado e aprendido muitos ou todos os conteúdos.
Na outra categoria, 72,5 \% optaram pelas mesmas categorias ( ver Gráfico 4.32 e Tabela II.16, no Anexo II ). 4.7 TEMPO GASTO PARA CONCLUIR A PROVA Ao responderem sobre o tempo de conclusão da prova ( Questão 9 ), quatro quintos dos estudantes ( 80,0 \% ) afirmaram ter gasto Entre duas e quatro horas ( Gráfico 4.33 a Gráfico 4.36 e, no Anexo II, as Tabelas II.17 e II.18 ).
Considerando -se as cinco Grandes Regiões brasileiras, os percentuais dos que utilizaram Entre duas e quatro horas para finalizar a prova, nas regiões Norte ( 75,4 \% ) e Nordeste ( 77,2 \% ), foram inferiores ao percentual nacional.
Nas demais Grandes Regiões, o percentual de alunos que dispensaram Entre duas e quatro horas para concluir a prova foi mais elevado que a média : 81,2 \%, na região Sudeste, 80,3 \%, na região Sul e 80,9 \%, na região Centro-Oeste, como mostrado no Gráfico 4.33.
A diferença entre as regiões não é estatisticamente significativa.
Uma vez analisadas as alternativas escolhidas pelos estudantes que se situam nos diferentes quartos de desempenho, em linhas gerais, observa -se uma tendência crescente : uma proporção maior de participantes, no quarto superior, declarou ter gasto Entre duas e quatro horas para concluir a prova quando comparada com a proporção de participantes nos quartos inferiores.
Há diferença estatisticamente significativa entre o primeiro e o último quartos de desempenho.
Na análise por Organização Acadêmica, os percentuais dos que utilizaram Entre duas e quatro horas para finalizar a prova, nas Universidades ( 77,9 \% ) e CEFET/IF ( 76,3 \% ), foram inferiores ao percentual nacional.
Nos demais tipos de Organização Acadêmica, o percentual de alunos que dispensaram Entre duas e quatro horas para concluir a prova foi maior do que a média : 83,8 \%, nas Faculdades e 81,5 \%, nos Centros Universitários.
Como mostrado no Gráfico 4.35, a diferença entre as organizações acadêmicas não é estatisticamente significativa.
Em relação às Categorias Administrativas, observa -se que, nas IES Públicas, 73,8 \% declararam ter gasto Entre duas e quatro horas para concluir a prova, sendo 80,7 \%, nas IES Privadas, com a mesma resposta.
A diferença entre os alunos que optaram por esse motivo de dificuldade nas duas categorias não é estatisticamente significativa ( ver Gráfico 4.36 e Tabela II.18, no Anexo II ).

% matched lemmas: 
\end{textsample}
